\documentclass{article}
\usepackage[utf8]{inputenc}
\usepackage[T1]{fontenc}
\usepackage{geometry}
\usepackage{graphicx}
\usepackage{float}
\usepackage{hyperref}
\usepackage[english,indonesian]{babel}

\title{Optimasi Pemrosesan Template LaTeX Menggunakan Pendekatan Berbasis AI}
\author{Academic Support Agent}
\date{\today}

\begin{document}
\maketitle

\begin{abstract}

\end{abstract}

\chapter{PENDAHULUAN}


\chapter{TINJAUAN PUSTAKA}


% Original content (commented out):
% \documentclass{article}
\usepackage[utf8]{inputenc}
%allow simultaneous greek and english input
\usepackage[greek,english]{babel}
\usepackage{alphabeta}
%package needed for image template
\usepackage{eso-pic}
%colour links, hrefs, etc whatever colour you like
%use only if needed. Shows warnings on Greek characters but nothing bad actually happens.
%just annoys you

%\usepackage[colorlinks = true,
            %linkcolor = black,
            %urlcolor  = black,
            %citecolor = blue,
            %anchorcolor = blue]{hyperref}

%Change default LaTeX contents name and figure name to Greek
\addto\captionsenglish{
    \renewcommand{\contentsname}{Περιεχόμενα}
    \renewcommand{\figurename}{Εικόνα}}
    %create a macro named BackgroundPic for setting a background on the first page
\newcommand\BackgroundPic{
    \put(-3,0){
    \parbox[b][\paperheight]{\paperwidth}{%
    \vfill
    \centering
    \includegraphics[width=\paperwidth,height=\paperheight]{LATEX_TEMPLATE_V2.pdf}
    \vfill
    }}}
%Info
\title{\vspace{0.2cm}\textbf{Τίτλος Εργασίας}}
\date{\textbf{ΗΜΕΡΟΜΗΝΙΑ\\ Version ΑΡΙΘΜΟΣ\\[8pt] ΟΝΟΜΑ ΜΑΘΗΜΑΤΟΣ}}
    \author{\textbf{AUTHOR HERE}}
%include page geometry packages
\usepackage{natbib}
\usepackage{graphicx}
\usepackage[a4paper, total={6.5in, 8in}]{geometry}
%document starts here, who would have imagined that?
\begin{document}
%set LATEX_TEMPLATE_V2 as template JUST for this page
%the asterisk means it will be set just for this page
\AddToShipoutPicture*{\BackgroundPic}
 \maketitle
\maketitle
%leave a bunch of empty space for the contents. Change cm if needed.
\vspace{4cm}
%shows contents automatically. Every section, subsection, etc automatically shows up here
\tableofcontents{}
%change page
\pagebreak
%create new command about pics on all other pages
\ClearShipoutPicture
\newcommand\Bgpic{
    \put(-4,0){
    \parbox[b][\paperheight]{\paperwidth}{%
    \vfill
    \centering
    \includegraphics[width=\paperwidth,height=\paperheight]{LATEX_TEMPLATE_V2.pdf}
    \vfill
    }}}
%set background picture for EVERY page after page 2, including page 2
%see, no asterisk
\AddToShipoutPicture{\Bgpic}
%yay, our first section, who would have thought

\section{Εισαγωγή}
\large
Lorem ipsum dolor sit amet, consectetur adipiscing elit. Aliquam scelerisque sit amet urna et interdum. Curabitur a scelerisque velit, quis accumsan dui. Sed risus mi, imperdiet quis nulla a, dictum vehicula nisl. Nulla vel erat aliquam, lobortis erat sit amet, faucibus ipsum. Sed tristique odio dolor, nec rutrum ante rutrum at. Aenean aliquam cursus mi, a condimentum nisi placerat non. Donec vel sem non elit consectetur sodales ut quis nibh. Cras ut ultrices leo. Etiam magna quam, accumsan sed erat ac, ultrices molestie nisi. Praesent ut lorem a magna maximus laoreet. Nulla non posuere enim.

Aenean placerat tellus id lorem semper, ut venenatis lorem accumsan. Integer varius, ligula quis tempor tristique, libero sem hendrerit tortor, sed semper sem nulla vel tortor. Integer in aliquam lacus. Aliquam at pretium lectus, et ornare odio. Sed at lorem a turpis euismod consectetur sed non risus. Cras sit amet nibh sit amet orci condimentum euismod. Donec placerat, felis vitae vehicula pulvinar, elit felis semper dui, ut interdum augue neque vulputate ipsum. Sed lacinia sodales ante id condimentum. Quisque sollicitudin nunc velit, at aliquet est interdum in.

Duis maximus sem quis mauris bibendum, id tempus nisi sodales. Maecenas a iaculis magna. Sed ac est elit. Vestibulum consequat magna quis orci finibus maximus. Duis malesuada rhoncus nunc non tristique. Aliquam ut interdum tortor, vel laoreet risus. In ultricies leo ut fringilla eleifend. Integer id metus a est imperdiet porta. Pellentesque habitant morbi tristique senectus et netus et malesuada fames ac turpis egestas. Vivamus ullamcorper, lacus at sagittis aliquet, orci arcu semper magna, id cursus felis lectus ut sem. Curabitur et dolor posuere, pulvinar erat eget, tempus lacus. Etiam id purus sed orci vestibulum aliquam. Sed suscipit ornare orci a tempus. Vestibulum ante lacus, volutpat vitae felis nec, cursus dictum elit.
	
\section{Κεφάλαιο Τάδε}

Donec posuere maximus massa in volutpat. Aliquam consectetur gravida massa sit amet egestas. Praesent facilisis congue urna sed maximus. Suspendisse potenti. Sed quis vehicula leo, vitae ultricies tellus. Cras a enim bibendum, vulputate metus quis, malesuada odio. Nam semper, tortor eget eleifend lacinia, felis massa auctor eros, sit amet aliquet quam ante sed libero. Donec at dapibus libero. Sed sem nunc, efficitur id efficitur laoreet, varius eu sapien. Morbi ullamcorper mi metus, ut tristique sem lobortis finibus. Mauris id feugiat magna, a condimentum nulla. Aliquam hendrerit consequat aliquet. Aliquam ullamcorper pellentesque nunc ac lobortis. Mauris non blandit lectus, in viverra ligula.

Nam non mattis lectus, ut ornare leo. Morbi mollis metus ut luctus dictum. Pellentesque quis neque at mauris porttitor maximus in eu urna. Fusce at ante in eros rhoncus tempus. Sed eu nisl ac enim vestibulum fermentum. Quisque lectus leo, eleifend id molestie id, tempor et mi. Nunc malesuada tellus nec arcu tristique, quis vehicula dui volutpat. Quisque in mattis elit. Integer id justo elementum, auctor dui sit amet, tristique turpis. Lorem ipsum dolor sit amet, consectetur adipiscing elit. Nunc pharetra et urna eget auctor. Phasellus accumsan arcu ut facilisis pretium. Quisque tincidunt felis a lacus ultrices, quis ornare libero commodo. Vivamus commodo magna non massa vehicula, eget consequat lacus commodo. Interdum et malesuada fames ac ante ipsum primis in faucibus.

\mbox{}
\vfill
\section{Βιβλιογραφία}
\small
[1] Συγγραφέας, Βιβλίο, Έτος\\[0pt]
[2] Συγγραφέας 2, Βιβλίο 2, Έτος 2\\[0pt]

\end{document}


\end{document}